\section{Related Work}

\subsection{Small Object Detection in Aerial Imagery}
Small object detection is fundamentally limited by low signal-to-noise ratios, reduced effective receptive fields, and aggressive downsampling that removes high-frequency cues.
In drone-captured imagery (e.g., traffic and crowd scenes), these issues are amplified by strong background clutter and heavy occlusion.
Common remedies include multi-scale feature pyramids, super-resolution or detail-preserving branches, and context-aware modules that leverage surrounding semantics to stabilize classification and localization.

\subsection{Efficient Feature Aggregation and Re-parameterization}
Dense feature aggregation designs (e.g., ELAN-style layer aggregation) improve gradient flow and preserve shallow spatial details that benefit small objects.
Re-parameterization further enables training-time multi-branch expressiveness (e.g., multi-kernel or identity branches) while collapsing to a single efficient convolution at inference, achieving high capacity without latency overhead.

\subsection{Large-kernel and Spatial Attention}
Large-kernel operators and spatial attention mechanisms expand the context window beyond local $3\\times3$ neighborhoods, which is particularly useful when small targets are ambiguous in isolation.
Recent works show that decomposed or depthwise large-kernel attention can provide long-range interactions at acceptable cost, often outperforming naïve kernel enlargement.

\subsection{Class Imbalance and Long-tail Reweighting}
Long-tail distributions are prevalent in aerial detection benchmarks.
Loss reweighting, focal-style objectives, sampling strategies, and adaptive priors are widely used to prevent dominant classes from overwhelming rare categories.
SC-ELAN complements backbone-level context modeling with training-time head reweighting (DetectCAI), aiming to improve tail-class recall and calibration without altering the inference contract.
